\chapter{Implementazione del parser}

Il contributo di questa tesi è l’implementazione di un tool per l’esecuzione di programmi scritti nel linguaggio \textit{Racing Choreographies}. Il sistema realizzato comprende due componenti principali: un front-end del linguaggio, responsabile dell’analisi lessicale e grammaticale dei programmi e della costruzione della loro rappresentazione interna (AST), e un simulatore che esegue i programmi secondo la semantica operazionale presentata nel Capitolo~2. 

L’architettura del tool è organizzata come una pipeline di trasformazioni successive che partono dal codice sorgente del programma e producono rappresentazioni via via più strutturate.

\begin{figure}[htbp]
\centering
\resizebox{0.78\textwidth}{!}{%
\begin{tikzpicture}[
    node distance=0.75cm,
    artifact/.style={
        draw,
        rounded corners,
        minimum width=4.2cm,
        minimum height=0.8cm,
        align=center
    },
    process/.style={
        draw,
        rectangle,
        minimum width=3.8cm,
        minimum height=0.75cm,
        align=center,
        fill=gray!10
    },
    arrow/.style={->, thick, >=Latex},
    groupbox/.style={
        draw,
        dashed,
        rounded corners,
        inner sep=0.3cm
    }
]

\node[artifact] (src) {File sorgente\\\texttt{.rc}};
\node[process, below=of src] (antlr) {Lexer + Parser\\ANTLR4};
\node[artifact, below=of antlr] (cst) {Concrete Syntax Tree\\(CST)};
\node[process, below=of cst] (builder) {AstBuilderVisitor};
\node[artifact, below=of builder] (ast) {Abstract Syntax Tree\\(AST)};
\node[process, below=of ast] (validator) {Validator};
\node[artifact, below=of validator] (vast) {Programma validato};
\node[process, below=of vast] (sim) {Simulatore};
\node[artifact, below=of sim] (exec) {Simulazione della coreografia};

\draw[arrow] (src) -- (antlr);
\draw[arrow] (antlr) -- (cst);
\draw[arrow] (cst) -- (builder);
\draw[arrow] (builder) -- (ast);
\draw[arrow] (ast) -- (validator);
\draw[arrow] (validator) -- (vast);
\draw[arrow] (vast) -- (sim);
\draw[arrow] (sim) -- (exec);

\node[groupbox, fit=(antlr)(cst)(builder)(ast)(validator)(vast),
      label=right:{Front-end}] {};

\end{tikzpicture}%
}
\caption{Pipeline generale del tool}
\label{fig:pipeline-tool}
\end{figure}

Nella prima fase il codice sorgente del programma viene passato in input a un \textit{lexer}, che produce una sequenza di tokens. Questa sequenza viene poi fornita in input a un \textit{parser}, il quale costruisce un \textit{Concrete Syntax Tree} (CST), ovvero una rappresentazione ad albero che riflette fedelmente la struttura sintattica del programma secondo la grammatica del linguaggio, includendo tutti i dettagli derivanti dalle regole grammaticali.

Successivamente, il CST viene trasformato in un \textit{Abstract Syntax Tree} (AST). L’AST rappresenta una versione più compatta e astratta della struttura sintattica del programma: elimina gli elementi sintattici non necessari (ad esempio nodi intermedi della grammatica o simboli ausiliari) e mantiene solo le informazioni rilevanti per il modello del linguaggio come descritto nel Capitolo~2.

Una volta costruito l’AST, il tool esegue una fase di validazione che verifica alcune proprietà strutturali dei programmi. Se il programma risulta valido, l’AST può essere utilizzato dal simulatore (Capitolo~5) per eseguire la coreografia secondo la semantica operazionale vista nel Capitolo~2.

In questo capitolo viene descritto il front-end del tool.

Il flusso generale del sistema è rappresentato nella Figura~\ref{fig:pipeline-tool}.

\section{Architettura del front-end}

Il front-end è costruito utilizzando il framework ANTLR4 \cite{antlr4tool}, che permette di generare automaticamente lexer e parser in diversi linguaggi di programmazione; per questo tool è stato utilizzato C++. A partire dalla grammatica definita nel Capitolo~3, ANTLR genera il codice necessario per riconoscere la struttura sintattica dei programmi e costruire il loro albero sintattico. L’implementazione completa del tool è disponibile nel repository del progetto \cite{rcparserrepo}. In particolare, lexer e parser si trovano nella directory \texttt{/generated}, mentre la grammatica è definita nel file \texttt{RacingChoreo.g4} all’interno della directory \texttt{/grammar}.

Il flusso del front-end generato da ANTLR è descritto nel seguente schema:

\begin{center}
\begin{tabular}{c}
\textbf{File sorgente} (\texttt{.rc}) \\
$\downarrow$ \\
\textbf{Lexer} (ANTLR) \\
$\downarrow$ \\
\textbf{Parser} (ANTLR) \\
$\downarrow$ \\
\textbf{Concrete Syntax Tree (CST)}
\end{tabular}
\end{center}

Supponiamo ora di avere un \textit{file.rc} scritto secondo la sintassi concreta vista nel Capitolo~2 che rappresenta il programma dell'esempio introdotto nei capitoli precedenti, in cui due worker $w1$ e $w2$ competono tramite una race per inviare un valore al processo $s$. Consideriamo ora questo frammento di codice del programma:

\begin{tcolorbox}[colback=white,colframe=black!35,boxrule=0.4pt]
\begin{lstlisting}
race s[k] : w1.r , w2.r -> s.ans;
\end{lstlisting}
\end{tcolorbox}
Quando questo programma viene dato in input al tool, la prima trasformazione consiste nella conversione del codice sorgente in una sequenza di tokens. Il tool permette di visualizzare questo risultato tramite il comando:

\begin{tcolorbox}[colback=black!2,colframe=black!35,boxrule=0.4pt]
\begin{verbatim}
> rc_parser tokens <file.rc>
\end{verbatim}
\end{tcolorbox}

L’output del lexer mostra i token riconosciuti nel programma insieme alla loro posizione nel file sorgente. Ad esempio, per l’istruzione precedente si ottiene:

\begin{center}
\end{center}

\begin{samepage}
\begin{verbatim}
6:4  RACE   "race"
6:9  ID     "s"
6:10 LBRACK "["
6:11 ID     "k"
6:12 RBRACK "]"
6:14 COLON  ":"
6:16 ID     "w1"
6:18 DOT    "."
6:19 ID     "r"
6:21 COMMA  ","
6:23 ID     "w2"
6:25 DOT    "."
6:26 ID     "r"
6:28 ARROW  "->"
6:31 ID     "s"
6:32 DOT    "."
6:33 ID     "ans"
6:36 SEMI   ";"
\end{verbatim}
\end{samepage}

Ogni riga dell’output indica la posizione del token nel file sorgente (linea e colonna), il tipo di token riconosciuto e il testo corrispondente nel programma.

La sequenza di token prodotta dal lexer costituisce l’input per il parser, che costruisce il \textit{Concrete Syntax Tree}. Il CST rappresenta la struttura sintattica del programma e contiene nodi che riflettono direttamente le produzioni grammaticali.

Nel caso dell’istruzione di race il parser costruisce un sottoalbero del CST che segue direttamente le produzioni della grammatica utilizzate per riconoscere tale istruzione. I nodi dell’albero corrispondono quindi sia ai token del programma (ad esempio \texttt{RACE}, \texttt{ID}, \texttt{COLON}, \texttt{ARROW}) sia alle regole grammaticali che strutturano l’istruzione.

Di conseguenza il CST contiene nodi per ciascun elemento sintattico della frase inclusi simboli puramente strutturali come parentesi quadre, virgole e separatori. Per l’istruzione di race l’albero include, ad esempio, nodi distinti per il costrutto \texttt{race}, per l’identificatore del processo ricevente \texttt{s}, per la variabile \texttt{k}, per le due espressioni concorrenti \texttt{w1.r} e \texttt{w2.r}, per il simbolo \texttt{->} e per la destinazione \texttt{s.ans}.

Questa rappresentazione è utile per verificare che il programma rispetti la sintassi del linguaggio, ma risulta troppo dettagliata per le fasi successive del sistema. Molti nodi dell’albero rappresentano infatti solo elementi della sintassi concreta e non corrispondono ai costrutti del modello del linguaggio. Per questo motivo il CST viene trasformato in un \textit{Abstract Syntax Tree} (AST).
\section{Abstract Syntax Tree}

Il front-end del tool introduce una seconda rappresentazione interna del programma ovvero l’\textit{Abstract Syntax Tree} (AST) che rappresenta il programma in una forma più compatta e più vicina alla struttura del linguaggio. In particolare, mentre il CST riflette il modo in cui il programma è stato riconosciuto dalla grammatica, l’AST mette direttamente in evidenza i costrutti del linguaggio, come procedure, blocchi di istruzioni, interazioni e costrutti di controllo.

Consideriamo ora il \textit{file.rc} del programma dell'esempio introdotto in precedenza per vedere come viene rappresentato. Il tool permette di stampare una rappresentazione testuale sul terminale dell’albero tramite il comando:

\begin{tcolorbox}[colback=black!2,colframe=black!35,boxrule=0.4pt]
\begin{verbatim}
> rc_parser ast file.rc
\end{verbatim}
\end{tcolorbox}
L’output relativo al programma di esempio è il seguente:
\begin{tcolorbox}[colback=black!2,colframe=black!35,boxrule=0.4pt]
\small
\begin{verbatim}
Program
  Procedures (1)
    ProcDef Request(c,w1,w2,s)
      Block (6 stmt)
        InteractionStmt
          Comm s.start -> w1.start
        InteractionStmt
          Comm s.start -> w2.start
        InteractionStmt
          Assign w1.r = 7
        InteractionStmt
          Assign w2.r = 5
        InteractionStmt
          Race s[k] : w1.r , w2.r -> s.ans
        IfRace (s[k])
        Then:
          Block (3 stmt)
            InteractionStmt
              Comm s.ans -> c.res
            InteractionStmt
              Select s -> c [FromW1]
            InteractionStmt
              Discharge s[k] : w2 -> s.lost
        Else:
          Block (3 stmt)
            InteractionStmt
              Comm s.ans -> c.res
            InteractionStmt
              Select s -> c [FromW2]
            InteractionStmt
              Discharge s[k] : w1 -> s.lost
  Main
    Block (1 stmt)
      Call Request(c,w1,w2,s)
\end{verbatim}
\end{tcolorbox}

L'implementazione completa dell'AST si trova nella directory \texttt{/src} nel repository del progetto \cite{rcparserrepo}. Di seguito vengono mostrate sono le principali parti di codice per descrivere le scelte progettuali che sono state fatte.

Alla radice dell’albero si trova il nodo \texttt{Program}, che contiene l’insieme delle definizioni di procedura e il blocco principale \texttt{Main}. Questa organizzazione riflette direttamente la struttura dei programmi del linguaggio.

In tutte le strutture dell’AST compare inoltre il campo \texttt{SourceRange loc}. Questo campo rappresenta la posizione del costrutto nel programma sorgente e viene utilizzato per associare ai nodi dell’albero le informazioni di posizione (file, riga e colonna). In questo modo il tool può produrre messaggi di errore che fanno riferimento direttamente al codice sorgente.

\begin{tcolorbox}[colback=white,colframe=black!35,boxrule=0.4pt]
\begin{lstlisting}
struct Program {
    std::vector<std::unique_ptr<ProcDef>> procedures;
    std::unique_ptr<Main> main;
    SourceRange loc;
};
\end{lstlisting}
\end{tcolorbox}

Nell’implementazione il nodo radice è definito come una struttura che contiene una lista di procedure e il nodo principale.

\begin{tcolorbox}[colback=white,colframe=black!35,boxrule=0.4pt]
\begin{lstlisting}
struct Block {
    std::vector<std::unique_ptr<Stmt>> statements;
    SourceRange loc;
};
\end{lstlisting}
\end{tcolorbox}

Si osserva che il corpo della procedura \texttt{Request} è rappresentato come un nodo \texttt{Block} contenente una sequenza di statement questo perchè, invece di mantenere nell’albero la struttura ricorsiva della grammatica, l’AST rappresenta esplicitamente i blocchi come sequenze di istruzioni.

\begin{tcolorbox}[colback=white,colframe=black!35,boxrule=0.4pt]
\begin{lstlisting}
using Stmt = std::variant<
    InteractionStmt,
    CallStmt,
    IfLocalStmt,
    IfRaceStmt
>;
\end{lstlisting}
\end{tcolorbox}

Gli statement del linguaggio comprendono diverse forme sintattiche, tra cui interazioni, chiamate di procedura e condizionali. Nell’implementazione queste alternative sono modellate tramite un tipo somma.

Questa soluzione consente di rappresentare in modo tipizzato le diverse forme che uno statement può assumere evitando di utilizzare gerarchie di classi.

L’AST non viene generato direttamente dal parser. Il parser di ANTLR produce infatti un CST che riflette la grammatica del linguaggio ma la costruzione dell’AST consiste in una trasformazione a partire dal CST.

Questa trasformazione è implementata tramite il meccanismo dei \textit{visitor} fornito da ANTLR. In particolare, la classe \texttt{AstBuilderVisitor} è un’estensione della classe \texttt{RacingChoreoBaseVisitor} generata automaticamente da ANTLR, nella quale vengono ridefiniti i metodi di visita necessari per costruire i nodi dell’AST.

Il tool definisce una classe \texttt{AstBuilderVisitor} che attraversa ricorsivamente il CST e costruisce i nodi dell’AST.

\begin{tcolorbox}[colback=white,colframe=black!35,boxrule=0.4pt]
\begin{lstlisting}
class AstBuilderVisitor : public RacingChoreoBaseVisitor {
public:
    std::any visitProgram(
        RacingChoreoParser::ProgramContext* ctx) override;
};
\end{lstlisting}
\end{tcolorbox}

Durante la visita dei nodi del CST, il visitor estrae le informazioni rilevanti e costruisce i nodi corrispondenti dell’AST.

Consideriamo nuovamente l’istruzione di \texttt{race} presente nel programma come esempio per capire il funzionamento:

\begin{tcolorbox}[colback=white,colframe=black!35,boxrule=0.4pt]
\begin{lstlisting}
race s[k] : w1.r , w2.r -> s.ans;
\end{lstlisting}
\end{tcolorbox}

Nel CST, come abbiamo visto in precedenza, questa istruzione è rappresentata da diversi nodi sintattici derivanti dalla grammatica. Il visitor raccoglie invece solo le informazioni di nostro interesse e costruisce il nodo dell’AST:

\begin{tcolorbox}[colback=white,colframe=black!35,boxrule=0.4pt]
\begin{lstlisting}
struct Race {
    RaceId id;
    ProcExpr left;
    ProcExpr right;
    ProcVar target;
    SourceRange loc;
};
\end{lstlisting}
\end{tcolorbox}

Questa struttura rappresenta esplicitamente i componenti semantici della \texttt{race}: l’identificatore della competizione, le due espressioni concorrenti e la variabile del processo ricevente.
\section{Validazione del programma}

Una volta costruito l’\textit{Abstract Syntax Tree}, il front-end del tool esegue una fase di validazione del programma. Questa fase ha lo scopo di verificare alcune proprietà strutturali del programma.

Come mostrato nella pipeline descritta all’inizio del capitolo, la validazione si colloca immediatamente dopo la costruzione dell’AST e prima dell’esecuzione della simulazione. Il validatore opera quindi su una rappresentazione del programma già indipendente dalla grammatica.

Il compito principale di questa fase è individuare errori che non riguardano la sintassi del programma, ma la sua struttura o l’uso scorretto dei costrutti del linguaggio. In particolare, il validatore verifica le seguenti proprietà:

\begin{itemize}

\item \textbf{Corretta definizione e invocazione delle procedure.}  
Il validatore verifica che ogni chiamata di procedura faccia riferimento a una procedura effettivamente definita nel programma e che non ci siano riferimenti a nomi di procedura inesistenti.

\item \textbf{Coerenza dei parametri nelle chiamate di procedura.}  
Per ogni istruzione \texttt{call}, viene controllato che il numero di processi passati come argomenti corrisponda al numero di parametri formali dichiarati nella definizione della procedura.

\item \textbf{Uso coerente dei processi nelle interazioni.}  
Il validatore verifica che i processi utilizzati nelle comunicazioni, nelle selezioni e negli altri costrutti del linguaggio siano tra quelli dichiarati nel contesto della procedura corrente.

\item \textbf{Uso consistente degli identificatori di race.}  
Per i costrutti \texttt{race}, \texttt{if (s[k])} e \texttt{discharge}, viene controllata la coerenza dell’identificatore della race e del processo associato, in modo da garantire che tali riferimenti siano utilizzati in modo consistente all’interno della coreografia.

\item \textbf{Presenza delle informazioni necessarie per l’esecuzione.}  
Il validatore verifica che il programma contenga gli elementi minimi necessari per l’esecuzione. In particolare controlla la presenza del blocco \texttt{main}, che il programma sia composto correttamente dal punto di vista strutturale (ovvero zero o più definizioni di procedura seguite dalla coreografia principale) e che le definizioni di procedura siano coerenti evitando duplicazioni di nomi tra procedure.
\end{itemize}

Questi controlli sono implementati nel modulo \texttt{Validation} del front-end, che attraversa l’AST e verifica le proprietà richieste. In caso di errore, il validatore produce un messaggio associato alla posizione corrispondente nel file sorgente sfruttando le informazioni di localizzazione propagate durante la costruzione dell’AST.

\section{Gestione degli errori}

Gli errori possono emergere in tre momenti distinti del processo:

\begin{itemize}
\item durante l'analisi lessicale e sintattica;
\item durante la validazione strutturale del programma;
\item durante l'esecuzione della simulazione.
\end{itemize}
Nel front-end della pipeline gli errori lessicali e sintattici vengono intercettati tramite un \textit{error listener} che estende il meccanismo di gestione degli errori fornito da ANTLR.
Esso permette di intercettare gli errori sintattici tramite l’interfaccia \texttt{BaseErrorListener}. Nel tool è stata definita una classe \texttt{ErrorListener} che estende questo meccanismo per raccogliere in modo strutturato le informazioni sugli errori rilevati durante il parsing.
L’implementazione si trova nei file \texttt{ErrorListener} nella directory \texttt{/src} \cite{rcparserrepo}

\begin{tcolorbox}[colback=white,colframe=black!35,boxrule=0.4pt]
\begin{lstlisting}
struct SyntaxError {
    std::string file;
    size_t line;
    size_t column;
    std::string message;
    std::string offendingText;
};
\end{lstlisting}
\end{tcolorbox}

La classe ha una struttura dati che registra per ciascun errore:

\begin{itemize}
\item il file sorgente;
\item la posizione dell’errore (linea e colonna);
\item il messaggio generato dal parser;
\item il token che ha causato l’errore.
\end{itemize}
Quando il parser rileva un errore, ANTLR invoca automaticamente il metodo \texttt{syntaxError}. Il metodo salva le informazioni relative all’errore all’interno di una lista mantenuta dall’oggetto listener.
\subsection*{Errori lessicali}

Gli errori lessicali vengono individuati dal lexer quando incontra sequenze di caratteri che non corrispondono ad alcun token definito nella grammatica.
Consideriamo un programma banale per mostrare il funzionamento:
\begin{tcolorbox}[colback=white,colframe=black!35,boxrule=0.4pt]
\begin{lstlisting}
main { @ }
\end{lstlisting}
\end{tcolorbox}
Il simbolo \texttt{@} non è definito nella grammatica del linguaggio e il lexer produce quindi un errore di riconoscimento del token e si ottiene:

\begin{tcolorbox}[colback=black!2,colframe=black!35,boxrule=0.4pt]
\begin{verbatim}
err_lex_01.rc:1:7: error: token recognition error at: '@'
  main { @ }
         ^
\end{verbatim}
\end{tcolorbox}
L’errore indica la posizione esatta del carattere non riconosciuto nel file sorgente.

\subsection*{Errori sintattici}

Gli errori sintattici vengono rilevati dal parser quando la sequenza di token prodotta dal lexer non soddisfa le regole della grammatica del linguaggio.
Consideriamo una istruzione di \texttt{race} in cui manca la virgola tra le due espressioni concorrenti:

\begin{tcolorbox}[colback=white,colframe=black!35,boxrule=0.4pt]
\begin{lstlisting}
race p[k] : p.x q.y -> p.z;
\end{lstlisting}
\end{tcolorbox}
Il parser segnala:

\begin{tcolorbox}[colback=black!2,colframe=black!35,boxrule=0.4pt]
\begin{verbatim}
err_01.rc:4:18: error: missing ',' at 'q'
    race p[k] : p.x q.y -> p.z;
                    ^
\end{verbatim}
\end{tcolorbox}
In modo analogo vengono segnalati errori dovuti a simboli mancanti o strutture incomplete, come ad esempio:

\begin{itemize}
\item espressioni mancanti in un assegnamento;
\item variabili incomplete in una comunicazione;
\item blocchi mancanti nei costrutti condizionali.
\end{itemize}

\subsection*{Errori strutturali}

Come abbiamo visto nella sezione precedente una volta completato il parsing e costruito l’AST il tool esegue una fase di validazione del programma. In questa fase vengono individuati errori che non riguardano la sintassi del programma ma la sua struttura logica.
Consideriamo ad esempio il seguente frammento di codice nel blocco per capire come vengono gestiti questi tipi di errori\texttt{main}:

\begin{tcolorbox}[colback=white,colframe=black!35,boxrule=0.4pt]
\begin{lstlisting}
main {
    call Request(c, w1, w2, s);
}
\end{lstlisting}
\end{tcolorbox}
In questo caso il programma tenta di invocare la procedura \texttt{Request}, che non è stata definita tra le procedure del programma. Durante la fase di validazione il tool rileva questa inconsistenza e produce il seguente messaggio di errore:

\begin{tcolorbox}[colback=black!2,colframe=black!35,boxrule=0.4pt]
\begin{verbatim}
ok_01.rc:20:4: error: call to undefined procedure 'Request'
      call Request(c, w1, w2, s);
      ^
\end{verbatim}
\end{tcolorbox}

\subsection*{Errori di esecuzione}

Un ultima categoria di errori è quella degli errori che possono emergere durante l’esecuzione della coreografia nel simulatore.
La gestione di questi errori è implementata all’interno della macchina di simulazione e viene descritta nel Capitolo~5.